\scanpage{11}

\uline{Variantes} \unsure{débarrasser|débarrassées} d'hypothèses noethériennes...

\uline{Variante} Pour un morphisme $Y\to X$ \uline{non ramifié}.

\uline{N. B.} Il faudrait également \uline{moduler} le th. précédent ....

\centerline{\rule{2.6cm}{0.3pt}}
\centerline{\rule{2cm}{0.3pt}}
\centerline{\rule{0.8cm}{0.3pt}}

\uline{Théorème} Soit $f:X\unsure{\longrightarrow|\text{---}}S$ morphisme de type fini, $S$ loc noeth, $E$ cohérent sur $X$. Soit $F$ sur $S$ le ss-foncteur de $S$
\[
F(S')=\begin{cases}
\emptyset & \text{si }F_{S'}\ \text{non plat }\inserted{\text{(resp. non fid plat)}}\text{ sur }S'\\
\{\emptyset\}&\text{sinon}
\end{cases}
\]
Pour que ce foncteur soit représentable, il f. et s. \unsure{qu il|qu'il} commute aux limites projectives adiques d'anneaux artiniens.

Vérification des divers \unsure{pts|pts.}

$1^\circ)$ Faisceau pour fpqc Clair

$2^\circ)$ Commute aux $\underrightarrow{\lim}$ d'anneaux : EGA IV 8 et \deleted{\unsure{IV}} 11
% slanted
\margin{$X$ type fini sur $S$, $F$ \unsure{noeth|noeth.}}

$3^\circ)$ Condition de commutation aux $\underleftarrow{\lim}$ d'artiniens

$4^\circ)$ $5^\circ)$ $6^\circ)$ Si $A\to A'$ injectif ($A'$ fini sur $A$ local artinien) alors $F(\deleted{\unsure{A'}}A)\to F(A')$ bijectif i.e $F(A')\ne\emptyset\Rightarrow F(A)\ne\emptyset$ i.e. $X_{A'}$ plat (fid plat) sur $A'$ $\Longrightarrow X_A$ plat (fid plat) sur $A$ \quad EGA IV 11

$7^\circ)$ \unsure{on a|\ldots} $F(T_0)\ne\emptyset$, $F(\deleted{\unsure{H}}U)\ne\emptyset\Longrightarrow F(T)\ne\emptyset$ i.e. $X_{T_0}$ plat et $X_U$ plat $\Longrightarrow X_T$ plat \quad EGA IV 11.
% archive sheet number: 10
